% ---------------------------------------------------
% Resumen (español) y Abstract (inglés)
% No dejar línea en blanco antes de \palabrasclave / \keywords
% ---------------------------------------------------

\begin{resumen}
La eficiencia hidráulica de las turbinas de impulso de baja caída depende de la geometría de la boquilla y del tipo de rodete. En este trabajo se analizan los datos experimentales de (Lakdawala, 2025), quienes ensayaron rodetes Pelton y Turgo de eje vertical con cangilones impresos en 3D en 80 corridas planificadas según un diseño central compuesto, variando el ángulo de la boquilla (90°–95°), su diámetro (10–14~mm), la distancia entre la boquilla y el rodete (40–60~mm) y el número de boquillas (una o dos). El objetivo es ajustar modelos de regresión que expliquen la eficiencia hidráulica en función de estas variables, verificar sus supuestos y comparar los resultados con el modelo cuadrático de superficie de respuesta reportado en el artículo original. El análisis descriptivo muestra una eficiencia media de 0.654 (desviación estándar de 0.024), mayor en el rodete Turgo y con dos boquillas, y una relación negativa marcada con el ángulo de la boquilla.
\palabrasclave{regresión lineal múltiple, diseño central compuesto, superficie de respuesta, turbina Pelton, turbina Turgo, eficiencia hidráulica}
\end{resumen}

\begin{abstract}
The hydraulic efficiency of low-head impulse turbines depends on the nozzle geometry and the runner type. This work analyzes the experimental data of (Lakdawala, 2025), who tested vertical-axis Pelton and Turgo runners with 3D-printed buckets in 80 runs planned under a central composite design, varying the nozzle angle (90°–95°), nozzle diameter (10–14~mm), nozzle standoff distance (40–60~mm) and number of nozzles (one or two). The aim is to fit regression models that explain hydraulic efficiency as a function of these variables, check their assumptions and compare the results with the quadratic response surface model reported in the original article. The descriptive analysis shows a mean efficiency of 0.654 (standard deviation 0.024), higher for the Turgo runner and with two nozzles, and a marked negative relationship with the nozzle angle.
\keywords{multiple linear regression, central composite design, response surface, Pelton turbine, Turgo turbine, hydraulic efficiency}
\end{abstract}

\pendiente{Completar el resumen y el abstract con los resultados del modelo final cuando estén listos.}

